\chapter{Métricas de evaluación y evolución diferencial}
\section{El problema que se quiere resolver}
Un grafo $G=(V,E)$ reúne vértices $V$ y aristas $E$ que indican qué pares están conectados. Un ciclo hamiltoniano recorre cada vértice exactamente una vez y vuelve al inicial por una arista existente. Si hay $n$ vértices, una ruta puede escribirse $R=(v_1,v_2,\ldots,v_n)$ y se entiende que después de $v_n$ se intenta regresar a $v_1$; por eso tiene $n$ conexiones, no $n-1$. La pregunta de decisión es si existe alguna ruta con esas propiedades, un problema clásico de la teoría de NP-completitud.\cite{de1}

Hay que separar camino hamiltoniano y ciclo hamiltoniano. Un camino visita todos los vértices sin repetirlos, pero no necesita regresar al punto de partida. Un ciclo exige además la última arista que une el final con el inicio. Esa condición de cierre es fácil de olvidar cuando la ruta se guarda como una lista, y por eso cualquier comprobación debe tratar la lista como circular.

El problema del viajante (TSP) agrega costos a las aristas y busca el ciclo más barato; aquí la pregunta es de existencia. Aun así, ambos están relacionados: los algoritmos para el TSP pueden adaptarse al problema hamiltoniano.\cite{de2} El conjunto FHCP Challenge Set reúne 1001 grafos de 66 a 9528 vértices; se sabe que todos tienen ciclo hamiltoniano y fueron diseñados para resistir a las heurísticas habituales.\cite{de2,de3} Su dificultad no viene solo del tamaño, sino de tener pocos ciclos y una estructura que impide descartar aristas con facilidad.\cite{de2} En grafos densos, en cambio, suele haber muchos ciclos y uno aparece pronto.\cite{de4} En el concurso asociado, el mejor resultado resolvió 985 grafos y ningún equipo resolvió los 1001,\cite{de3} lo que da una idea de la escala del reto. El ejemplo pequeño de la sección~\ref{de:sec:ejemplo} fue inventado para explicar el método y no pertenece a ese conjunto.

Este trabajo tiene dos partes unidas por una misma preocupación. Como una búsqueda con azar no se comporta igual en cada corrida, primero hay que saber cómo medirla; después se estudia qué método podría usarse, la evolución diferencial, y qué exige adaptarlo.

\section{Métricas de evaluación}
\subsection{Qué es una métrica y qué se necesita para medir}
Una métrica es una regla que convierte el comportamiento de un algoritmo en un número comparable. Antes de calcular cualquiera conviene fijar el vocabulario. Una instancia es un caso concreto del problema, como un grafo. Una ejecución es una corrida completa del algoritmo sobre una instancia; los métodos con azar dan resultados distintos en cada una, y la semilla es el valor que fija esa secuencia de azar para poder repetirla. Una evaluación es calcular una vez la calidad de una solución candidata. El presupuesto es el máximo de recursos permitido a una ejecución.

Una sola corrida no describe un método estocástico: puede haber tenido suerte o mala fortuna. Por eso las medidas se resumen sobre varias ejecuciones independientes.\cite{de5} En general interesan tres preguntas distintas, que no conviene mezclar: ¿qué tan seguido lo logra? (fiabilidad), ¿cuánto le cuesta? (esfuerzo) y ¿qué tan cerca queda cuando no lo logra? (calidad). Eiben y Smith separan de forma parecida la velocidad, o eficiencia, de la calidad, o efectividad.\cite{de7}

\subsection{Tiempo, costo computacional, iteraciones y evaluaciones}
Las cuatro medidas parecen decir lo mismo, pero miden cosas diferentes.

\begin{center}

\begin{adjustbox}{max width=\linewidth}
\begin{tabular}{@{}P{2.6cm}P{6.6cm}P{6.6cm}@{}}
\toprule
Medida & Qué mide & Limitación\\
\midrule
Tiempo & Tiempo transcurrido o tiempo de procesador (CPU) usado por el programa. & Depende de la máquina, el lenguaje y la implementación.\\
Costo computacional & Recursos que consume el método: operaciones, tiempo o memoria. & Es un término amplio: hay que declarar qué recurso se cuenta y en qué unidad.\\
Iteraciones & Ciclos del algoritmo; en métodos evolutivos se llaman generaciones. & Una generación cuesta distinto según el tamaño de la población.\\
Evaluaciones & Veces que se calculó la calidad de una solución. & No todas cuestan lo mismo.\\
\bottomrule
\end{tabular}
\end{adjustbox}
\end{center}

Las evaluaciones son la medida independiente de la máquina más común en optimización continua, y la plataforma COCO las usa precisamente porque el tiempo de CPU depende del lenguaje, el estilo de programación y el equipo.\cite{de5,de6} La recomendación práctica es reportar evaluaciones y tiempo de CPU, y si se dan generaciones, incluir también las evaluaciones.\cite{de5} La advertencia es similar en Eiben y Smith: no conviene contar generaciones cuando se comparan configuraciones con poblaciones distintas.\cite{de7}

Ejemplo de la diferencia. En evolución diferencial cada generación evalúa una candidata por cada miembro de la población, además de la evaluación inicial. Con 20 individuos y 50 generaciones hay $20+20\cdot 50=1020$ evaluaciones; con 40 individuos y las mismas 50 generaciones hay $2040$. Decir ``50 generaciones'' oculta que el segundo caso costó el doble.

Propuesta para este problema: cada evaluación de una ruta revisa $n$ conexiones, de modo que evaluar en un grafo de 3000 vértices cuesta mucho más que en uno de 66. Por eso las evaluaciones solo se comparan directamente entre instancias de tamaño parecido, y el tiempo debe acompañarlas cuando el método añade trabajo interno.

\subsection{Tasa de éxito y ejecuciones fallidas}
Hablar de éxito exige definir un objetivo. La tasa de éxito es el porcentaje de ejecuciones que lo alcanzan dentro del presupuesto. A partir de ahí hay dos formas de plantear una medición: con objetivo fijo se pregunta cuánto esfuerzo costó llegar, y con presupuesto fijo se pregunta a qué calidad se llegó con ese esfuerzo. Hansen y colaboradores prefieren la primera porque da resultados cuantitativos e interpretables.\cite{de6}

El punto delicado son los fallos. En una ejecución fallida el tiempo hasta el éxito no está definido: solo se sabe cuánto consumió.\cite{de6} Quien mide con objetivo fijo debe declarar cómo los trata.\cite{de5} Si simplemente se eliminan, el método parece mejor de lo que fue. Las medidas habituales que sí los incluyen son:
\begin{itemize}[leftmargin=*]
\item aRT (average runtime, también llamado ERT): suma las evaluaciones de todas las ejecuciones y divide entre el número de éxitos. Estima cuánto costaría llegar a un éxito si se reiniciara el algoritmo cada vez que falla.\cite{de5,de6}
\item PAR-$k$: promedia lo consumido contando cada fallo como $k$ veces el presupuesto máximo; PAR2 y PAR10 son los casos comunes.\cite{de5}
\item Mediana con fallos al final: los fallos se ordenan por encima de todo éxito. Sigue siendo válida mientras al menos la mitad de las ejecuciones tenga éxito, y los cuantiles son más robustos que la media.\cite{de5}
\end{itemize}

Ejemplo con diez ejecuciones inventadas. Con un presupuesto de 1000 evaluaciones, siete ejecuciones tuvieron éxito con 120, 150, 180, 200, 240, 300 y 410 evaluaciones, y tres fallaron.

\begin{center}

\begin{adjustbox}{max width=\linewidth}
\begin{tabular}{@{}P{4.8cm}P{3.5cm}rP{5.3cm}@{}}
\toprule
Resumen & Cálculo & Valor & Qué muestra\\
\midrule
Tasa de éxito & $7/10$ & 70\,\% & Fiabilidad.\\
Media solo de éxitos & $1600/7$ & 228.6 & Oculta los tres fallos.\\
Mediana con fallos al final & entre 240 y 300 & 270 & Una ejecución típica.\\
aRT & $(1600+3000)/7$ & 657.1 & Esfuerzo total por cada éxito.\\
PAR10 & $(1600+30\,000)/10$ & 3160 & Castiga con fuerza los fallos.\\
\bottomrule
\end{tabular}
\end{adjustbox}
\end{center}

La misma información da 228.6 o 657.1 según cómo se trate a los fallos, y por eso el informe debe decir qué medida usó. Además, una tasa calculada con pocas ejecuciones es incierta. Para muestras pequeñas se recomiendan los intervalos de Wilson o Jeffreys, porque el intervalo clásico de Wald tiene una cobertura errática.\cite{de8} En este ejemplo, $7/10$ da un intervalo de Wilson al 95\,\% de aproximadamente 40\,\% a 89\,\%: el ``70\,\%'' es mucho menos exacto de lo que aparenta.

Propuesta: registrar la causa de cada fallo (presupuesto agotado, límite de tiempo o error técnico), porque un error del programa no informa sobre el método. Y si ninguna ejecución tiene éxito, el aRT no está definido, así que debe reportarse ``0 de $N$'' junto con el presupuesto usado.

\subsection{Precisión, recall y F1}
Estas medidas vienen de la clasificación, donde cada elemento se etiqueta como positivo o negativo. Un verdadero positivo ($TP$) se señaló como positivo y lo era; un falso positivo ($FP$) se señaló sin serlo; un falso negativo ($FN$) era positivo y no se señaló. Con esos conteos:\cite{de9}
\begin{equation}
\text{Precisión}=\frac{TP}{TP+FP},\qquad
\text{Recall}=\frac{TP}{TP+FN},\qquad
F_1=\frac{2\,TP}{2\,TP+FP+FN}.\label{de:eq:prf}
\end{equation}
La precisión responde: de lo que señalé, ¿cuánto era correcto? El recall o exhaustividad responde: de lo que debía encontrar, ¿cuánto encontré? F1 es su media armónica. Ninguna usa los verdaderos negativos, y pares distintos de precisión y recall pueden dar el mismo F1, por lo que no debe reportarse solo.\cite{de9}

Ejemplo de clasificación. De 12 mensajes de spam, un filtro marcó 8 correctamente y además marcó 2 mensajes normales. Entonces $TP=8$, $FP=2$ y $FN=4$. La precisión es $8/10=0.80$, el recall es $8/12\approx 0.67$ y F1 es aproximadamente $0.73$.

Lo importante es que estas fórmulas no significan nada hasta responder cuatro preguntas: (1) qué elemento se clasifica, (2) qué significa positivo, (3) cuál es la referencia que dice qué es correcto y (4) en qué nivel se mide: conexión, ruta, ejecución o instancia. En el spam las respuestas son obvias; en el ciclo hamiltoniano no.

\subsection{Por qué no basta con un porcentaje de conexiones válidas}
Las fuentes revisadas no ofrecen una definición estándar de estas medidas para el ciclo hamiltoniano. La siguiente tabla es una propuesta que aplica la ecuación~\eqref{de:eq:prf} al grafo de la sección~\ref{de:sec:ejemplo}, que tiene 9 aristas. $M$ es la ruta inválida $(A,D,C,B,F,E)$, a la que solo le falta la conexión de regreso $EA$. $C_3$ es un ciclo perfecto y $C_1$ es otro ciclo del mismo grafo.

\begin{center}

\begin{adjustbox}{max width=\linewidth}
\begin{tabular}{@{}P{4.9cm}ccP{6.5cm}@{}}
\toprule
Qué se toma como positivo & Caso & P / R / F1 & Problema\\
\midrule
La conexión existe en el grafo & $M$ & .83 / .56 / .67 & La precisión es el porcentaje de conexiones válidas.\\
 & $C_3$ & 1 / .67 / .80 & Un ciclo perfecto nunca llega a recall 1: usa 6 de las 9 aristas.\\
\addlinespace
La conexión pertenece a un ciclo de referencia ($C_1$) & $M$ & .50 / .50 / .50 & Las tres medidas coinciden.\\
 & $C_3$ & .50 / .50 / .50 & Un ciclo perfecto empata con una ruta inválida.\\
\addlinespace
La ejecución entrega un ciclo verificado & 10 ejec. & 1 / .70 / .82 & Solo repite la tasa de éxito.\\
\bottomrule
\end{tabular}
\end{adjustbox}
\end{center}

Con un ciclo de referencia, el resultado no es casualidad. La ruta y la referencia tienen $n$ conexiones cada una, así que $TP+FP=TP+FN=n$, y por tanto precisión, recall y F1 valen siempre $TP/n$. Son tres nombres para un solo porcentaje, justo lo que el enunciado advierte que no basta. Peor aún, si el grafo tiene varios ciclos, la medida castiga una solución correcta solo por ser distinta de la elegida como referencia.

A nivel de ejecución tampoco se gana información. Si el éxito se declara después de verificar el ciclo, no hay falsos positivos y la precisión vale 1; el recall coincide con la tasa de éxito y $F_1=2R/(1+R)$ es solo una transformación suya. Además, como todos los grafos del FHCP tienen ciclo,\cite{de3} no existe clase negativa a nivel de instancia. Por eso esta investigación no adopta una definición: la deja planteada como duda para el profesor.

\subsection{Qué medidas parecen aplicables}
La siguiente valoración es una propuesta pendiente de confirmar.

\begin{center}

\begin{adjustbox}{max width=\linewidth}
\begin{tabular}{@{}P{4.6cm}P{7.6cm}P{3.4cm}@{}}
\toprule
Medida & Significado en el ciclo hamiltoniano & Valoración\\
\midrule
Tasa de éxito con intervalo & Porcentaje de ejecuciones con ciclo verificado. & Aplicable.\\
Evaluaciones y tiempo & Esfuerzo hasta el éxito o hasta agotar el presupuesto. & Aplicable.\\
aRT, PAR-$k$, mediana & Resumen de esfuerzo que incluye los fallos. & Transferible; confirmar.\\
Conexiones faltantes & Las que no existen en la mejor ruta final. & Calidad parcial.\\
Porcentaje de conexiones válidas & Avance parcial; no significa haber resuelto. & Con su propio nombre.\\
Precisión, recall y F1 & Dependen de la definición elegida. & Duda abierta.\\
\bottomrule
\end{tabular}
\end{adjustbox}
\end{center}

\section{Evolución diferencial: el algoritmo original}
\subsection{Origen e idea central}
La evolución diferencial (DE, Differential Evolution) fue propuesta por Storn y Price para optimizar funciones sobre valores continuos.\cite{de10} Sigue siendo una de las familias de optimizadores evolutivos más estudiadas, con una extensa literatura sobre variantes y parámetros.\cite{de11} Mantiene una población de $NP$ soluciones candidatas, cada una representada como un vector de $D$ números reales, y una función objetivo les asigna un costo que se quiere minimizar.

Su idea distintiva es simple: en lugar de usar una distribución de azar externa para decidir cuánto moverse, usa la diferencia entre dos soluciones de la propia población. Cuando la población está dispersa, las diferencias son grandes y los saltos también; cuando se concentra, los saltos se acortan solos. Frente al algoritmo genético, donde se seleccionan padres y se forman descendientes, en DE cada candidata compite directamente contra el vector al que quiere reemplazar.\cite{de10}

\subsection{Qué ocurre en una generación}
La variante básica se llama DE/rand/1/bin: el vector base se elige al azar (rand), se usa una diferencia (1) y el cruce es binomial (bin).\cite{de10} En cada generación, cada vector objetivo $x_i$ pasa por tres pasos:
\begin{enumerate}[leftmargin=*]
\item Mutación. Se eligen tres vectores distintos entre sí y distintos de $x_i$, y se forma el vector mutante
\begin{equation}
v=x_{r_1}+F\,(x_{r_2}-x_{r_3}),\label{de:eq:mut}
\end{equation}
donde $F$ es el factor de escala. Por eso se requieren al menos cuatro vectores.\cite{de10}
\item Cruce. Se forma el vector de prueba $u$. Cada componente viene de $v$ con probabilidad $CR$ y de $x_i$ en caso contrario, excepto una posición elegida al azar que siempre viene de $v$; así se garantiza tomar al menos una componente del donante, aunque el valor o la permutación decodificada pueden coincidir con los anteriores.\cite{de10}
\item Selección. Se evalúa $u$. Si cuesta menos que $x_i$, lo reemplaza en la siguiente generación; si no, $x_i$ se conserva.\cite{de10}
\end{enumerate}

Ejemplo numérico. Se quiere minimizar $f(x,y)=x^2+y^2$ con $x_1=(3,1)$, $x_2=(0,3)$, $x_3=(1,-1)$, $x_4=(-2,2)$, $F=0.5$ y $CR=0.5$. Para el objetivo $x_1$ se eligen $x_3$ como base y $x_2-x_4=(2,1)$ como diferencia, así que $v=(1,-1)+0.5\,(2,1)=(2,-0.5)$. En el cruce, la posición 2 es la forzada; en la posición 1 el azar da 0.7, mayor que $CR$, y se conserva el 3 de $x_1$. Queda $u=(3,-0.5)$, con costo $9.25$, menor que el costo 10 de $x_1$. Por eso $u$ ocupa su lugar en la siguiente generación.

\subsection{Parámetros que influyen}
Storn y Price sugieren $NP$ entre $5D$ y $10D$, nunca menos de 4; $F=0.5$ como punto de partida, con valores por debajo de 0.4 o por encima de 1 eficaces solo en ocasiones; y $CR=0.1$ como primera opción, probando primero 0.9 o 1 para ver si se obtiene una solución rápida. Si la población converge prematuramente, aconsejan aumentar $F$ o $NP$.\cite{de10} La tabla siguiente resume el efecto esperado a partir de esas indicaciones y de la ecuación~\eqref{de:eq:mut}.

\begin{center}

\begin{adjustbox}{max width=\linewidth}
\begin{tabular}{@{}P{1.5cm}P{7.1cm}P{7.1cm}@{}}
\toprule
Ajuste & Si se reduce & Si se aumenta\\
\midrule
$NP$ & Menos diversidad; cada generación cuesta menos. & Más diversidad y más evaluaciones por generación.\\
$F$ & Saltos cortos; riesgo de convergencia prematura. & Saltos largos; más exploración y búsqueda menos fina.\\
$CR$ & El vector de prueba se parece más al objetivo. & Hereda más del mutante; a menudo acelera la convergencia.\\
\bottomrule
\end{tabular}
\end{adjustbox}
\end{center}

Propuesta: esas recomendaciones son de 1997 y para problemas continuos, así que aquí solo sirven como punto de partida. Aplicadas a una ruta de 3000 vértices darían entre 15\,000 y 30\,000 vectores, algo que habría que justificar. La adaptación automática de parámetros es un tema desarrollado después.\cite{de11}

\subsection{Cuándo termina la búsqueda}
Los criterios de paro para DE se agrupan en cinco clases: por referencia (se alcanza un valor), por agotamiento (se acaba el presupuesto), por mejora (deja de mejorar), por movimiento (la población deja de moverse) y por distribución (la población se concentra).\cite{de12} En ese estudio, hecho sobre un problema real concreto, algunos criterios de mejora resultaron poco fiables, y los autores advierten que los resultados pueden cambiar en otros casos.\cite{de12}

\section{Adaptación a órdenes de visita}
\subsection{Por qué no se puede aplicar directamente}
DE necesita sumar, restar y escalar vectores, y una permutación, un orden de todos los vértices sin repetir, no admite esas operaciones con sentido. Si se numeran los vértices, $(1,2,3,4,5)$ y $(2,3,4,5,1)$ son el mismo ciclo con otro inicio, pero su resta es $(-1,-1,-1,-1,4)$. Aplicar la mitad de esa diferencia a $(1,3,2,4,5)$ produce $(0.5,\,2.5,\,1.5,\,3.5,\,7)$, que no es una ruta: los valores ni siquiera son vértices. La ecuación~\eqref{de:eq:mut} supone que estar ``cerca'' en los números significa parecerse como solución, y eso aquí no ocurre.

Hay además una redundancia propia del problema. En un grafo no dirigido, cada ciclo se escribe de $2n$ formas: $n$ puntos de inicio por 2 sentidos. Con 5 vértices, las 120 permutaciones representan solo 12 ciclos candidatos distintos. Fijar el vértice inicial elimina una parte de esa redundancia, y la del sentido queda.

\subsection{Función de evaluación y validez}
Sea $R=(v_1,\ldots,v_n)$ una permutación y $v_{n+1}=v_1$. Una evaluación simple, propuesta para este documento, cuenta las conexiones inexistentes:
\begin{equation}
b(R)=\sum_{i=1}^{n}\mathbf{1}\bigl[(v_i,v_{i+1})\notin E\bigr].\label{de:eq:b}
\end{equation}
La expresión $\mathbf{1}[P]$ vale 1 si $P$ es verdadera y 0 si no, de modo que $b(R)$ revisa exactamente $n$ conexiones, incluida la de regreso. Como DE minimiza, $b$ puede usarse directamente como costo: no hace falta transformarlo en un fitness que se maximice. En este documento la letra $F$ se reserva para el factor de escala de DE; conviene no confundirla con la $F(R)$ que otros métodos del equipo usen como fitness.

La validez se comprueba aparte: (1) la ruta contiene exactamente los $n$ vértices distintos; (2) cada par consecutivo es una arista; y (3) también existe la arista del último al primero. Si la representación garantiza una permutación, $b(R)=0$ equivale a validez, pero conviene conservar un verificador independiente para detectar errores. Una ruta con $b=1$ puede estar cerca según esta medida, pero nunca es una solución.

\subsection{Pseudocódigo de la adaptación}\label{de:sec:pseudo}
El esquema siguiente es deliberadamente general: conserva las fases de DE y deja indicados los puntos que cambian según la opción de la sección~\ref{de:sec:opciones}. No es un programa completo ni supone un lenguaje.
\begin{enumerate}[leftmargin=*,label=\arabic*.]
\item Recibir el grafo $G=(V,E)$, los parámetros $NP$, $F$ y $CR$, el presupuesto de evaluaciones $B$ y una representación de las rutas.
\item Crear $NP$ soluciones iniciales al azar, obtener su ruta y calcular $b$ con la ecuación~\eqref{de:eq:b}. Poner el contador de evaluaciones en $NP$ y guardar la mejor.
\item Si alguna tiene $b=0$, pasar al paso 10. Exigir $B\geq NP$ y pasar también al paso 10 si se agotó el presupuesto tras inicializar. Exigir $B\geq NP$ y pasar también al paso 10 si se agotó el presupuesto tras inicializar.
\item Para cada solución $x_i$, elegir tres soluciones distintas entre sí y distintas de $x_i$.
\item Mutar con ellas para obtener $v$. Ajuste A: con vectores es la ecuación~\eqref{de:eq:mut}; con permutaciones hay que definir la diferencia y la escala.
\item Cruzar $x_i$ y $v$ con probabilidad $CR$ para obtener $u$. Ajuste B: con permutaciones debe evitar vértices repetidos.
\item Obtener la ruta de $u$. Ajuste C: no existe en DE original; con claves consiste en ordenarlas.
\item Calcular $b$ y sumar uno al contador. Seleccionar: si la candidata mejora a $x_i$, ocupa su lugar en la siguiente generación; si no, $x_i$ se conserva. Ajuste D: decidir si un empate también se acepta. Actualizar la mejor ruta y, si $b=0$ o se agotó $B$, pasar al paso 10.
\item Al terminar la población, reemplazarla por la nueva y repetir desde el paso 4.
\item Comprobar la mejor ruta con el verificador independiente. Si es válida, devolverla junto con las evaluaciones usadas. Si no, informar ``no encontrado con este presupuesto de búsqueda''.
\end{enumerate}

Propuesta de paro: aquí existe un criterio de referencia natural, $b=0$, porque esa ruta se comprueba como ciclo. Debe combinarse con un presupuesto para que las ejecuciones fallidas también terminen. Reiniciar tras un estancamiento es otra posibilidad por discutir.

\section{Ejemplo didáctico de seis vértices}\label{de:sec:ejemplo}
Tomemos $V=\{A,B,C,D,E,F\}$ y el siguiente grafo no dirigido, definido por su lista completa de aristas:
\[
E=\{AB,BC,CD,DE,EF,FA,AD,BF,CE\}.
\]
Tiene tres ciclos hamiltonianos: $C_1$, $A\text{-}B\text{-}C\text{-}D\text{-}E\text{-}F\text{-}A$; $C_2$, $A\text{-}B\text{-}F\text{-}E\text{-}C\text{-}D\text{-}A$; y $C_3$, $A\text{-}D\text{-}E\text{-}C\text{-}B\text{-}F\text{-}A$. El ejemplo usa claves aleatorias, descritas en la sección~\ref{de:sec:opciones}, solo porque hacen visible la aritmética de DE; no significa que se hayan elegido como adaptación. Cada vector tiene una clave por vértice en el orden $A,B,C,D,E,F$. Para obtener la ruta, las claves se ordenan de menor a mayor y se leen los vértices correspondientes. Se usa $NP=4$, $F=0.5$ y $CR=0.5$.

\begin{center}

\begin{adjustbox}{max width=\linewidth}
\begin{tabular}{@{}lP{5cm}P{4.3cm}P{3cm}c@{}}
\toprule
Vector & Claves & Ruta (con regreso) & Ausentes & $b$\\
\midrule
$x_1$ & $(.3,\ .7,\ .1,\ .8,\ .9,\ .6)$ & $C\text{-}A\text{-}F\text{-}B\text{-}D\text{-}E\text{-}C$ & $CA,BD$ & 2\\
$x_2$ & $(.4,\ .2,\ .9,\ .7,\ .6,\ .8)$ & $B\text{-}A\text{-}E\text{-}D\text{-}F\text{-}C\text{-}B$ & $AE,DF,FC$ & 3\\
$x_3$ & $(.4,\ .3,\ .9,\ .7,\ .5,\ .2)$ & $F\text{-}B\text{-}A\text{-}E\text{-}D\text{-}C\text{-}F$ & $AE,CF$ & 2\\
$x_4$ & $(.9,\ .2,\ .7,\ .3,\ .6,\ .5)$ & $B\text{-}D\text{-}F\text{-}E\text{-}C\text{-}A\text{-}B$ & $BD,DF,CA$ & 3\\
\bottomrule
\end{tabular}
\end{adjustbox}
\end{center}

\noindent Paso 1: evaluar. Las cuatro rutas visitan los seis vértices, pero ninguna es válida. En $x_1$, la clave menor es 0.1 y pertenece a $C$, por lo que la ruta empieza en $C$. Evaluar la población costó 4 evaluaciones.

\noindent Paso 2: primer intento, rechazado. Para el objetivo $x_1$ se usa $v=x_2+0.5\,(x_3-x_4)=(.15,\ .25,\ 1.0,\ .9,\ .55,\ .65)$. Los azares por posición son .2, .8, .4, .9, .6 y .3, y la posición forzada es la 5. Se toman de $v$ las posiciones 1, 3, 5 y 6, lo que da $u=(.15,\ .7,\ 1.0,\ .8,\ .55,\ .65)$. Su ruta es $A\text{-}E\text{-}F\text{-}B\text{-}D\text{-}C\text{-}A$, y faltan $AE$, $BD$ y $CA$, así que $b=3$. Como 3 no es menor que 2, se rechaza y $x_1$ se conserva. Van 5 evaluaciones.

\noindent Paso 3: segundo intento, aceptado. Para el objetivo $x_2$ se usa $v=x_3+0.5\,(x_4-x_1)=(.7,\ .05,\ 1.2,\ .45,\ .35,\ .15)$. Los azares son .7, .1, .9, .3, .6 y .8, y la posición forzada es la 6. Se toman de $v$ las posiciones 2, 4 y 6, lo que da $u=(.4,\ .05,\ .9,\ .45,\ .6,\ .15)$. Su ruta es $R'=B\text{-}F\text{-}A\text{-}D\text{-}E\text{-}C\text{-}B$ con $b=0$. Como $0<3$, reemplaza a $x_2$ y la búsqueda se detiene con 6 evaluaciones.

\noindent Paso 4: comprobación independiente. Los vértices antes del regreso son $B,F,A,D,E,C$, todos distintos y todos los de $V$. La última columna de la tabla corresponde a la conexión de cierre. $R'$ es el ciclo $C_3$ escrito desde otro inicio.

\begin{center}

\begin{adjustbox}{max width=\linewidth}
\begin{tabular}{ccccccc}
\toprule
Posición & 1 & 2 & 3 & 4 & 5 & 6\\
\midrule
Conexión & $B\!\to\!F$ & $F\!\to\!A$ & $A\!\to\!D$ & $D\!\to\!E$ & $E\!\to\!C$ & $C\!\to\!B$\\
¿Existe? & Sí & Sí & Sí & Sí & Sí & Sí\\
\bottomrule
\end{tabular}
\end{adjustbox}
\end{center}

\noindent Contraejemplo del cierre. La ruta $M=(A,D,C,B,F,E)$ usa cinco conexiones existentes: $AD$, $DC$, $CB$, $BF$ y $FE$. Solo falla al regresar, porque $EA$ no es arista, así que $b=1$. Visita todos los vértices y aun así no es un ciclo: es un camino hamiltoniano.

Carácter del ejemplo. Las claves y los números al azar se escogieron para que el ejemplo fuera breve. Es una secuencia ilustrativa posible, no una ejecución medida ni un resultado del FHCP; hallar un ciclo en dos intentos no es lo esperable.

\subsection{Qué enseña el ejemplo}
El primer intento muestra que una candidata puede empeorar, y que la selección impide perder lo ganado. El segundo muestra que cambiar pocas claves puede reordenar toda la ruta: aunque $u$ conserva tres claves de $x_2$, su ruta solo conserva dos de las seis conexiones. Esa es una virtud para explorar y una dificultad para mejorar con precisión. Por último, $x_1$ y $x_3$ tenían $b=2$ y empataban sin parecerse como rutas, lo que anticipa el problema de las mesetas.

\section{Aplicación al ciclo hamiltoniano: opciones y dificultades}\label{de:sec:opciones}
Esta investigación no elige una opción. La decisión depende de las dudas de la sección~\ref{de:sec:dudas} y de pruebas que aquí no se realizaron.

Opción A: claves aleatorias. Se conserva DE continua y cada vector se ordena para leer una ruta; la técnica procede de los algoritmos genéticos\cite{de13} y se ha estudiado en DE para permutaciones, incluido el problema del viajante.\cite{de14} La regla del valor de posición menor (SPV) sigue una idea relacionada.\cite{de16,de17} Su ventaja es que toda clave da una permutación válida y DE no cambia. Su dificultad es que infinitos vectores representan la misma ruta, lo que crea mesetas en la búsqueda; Santucci y colaboradores señalan esto como una causa probable de su bajo rendimiento.\cite{de15}

Opción B: DE algebraica. Se trabaja directamente con permutaciones y se redefinen las operaciones: la diferencia entre dos rutas es una secuencia de intercambios adyacentes, y escalarla por $F$ significa aplicar solo una parte de ella.\cite{de15} Su ventaja es que no hay decodificación ni distinción entre vector y ruta. Su dificultad es que es más compleja y fue probada en secuenciación de tareas, no en este problema.

Opción C: otros enfoques discretos. El manual de Onwubolu y Davendra reúne enfoques como la transformación directa e inversa y los índices de posición relativa.\cite{de16} Queda pendiente revisarlos, porque solo se consultó su índice. Santucci y colaboradores también mencionan una DE geométrica aplicada al TSP con resultados pobres.\cite{de15}

Vacío detectado. Una búsqueda dirigida no localizó un artículo que aplique DE específicamente al ciclo hamiltoniano. Lo más cercano es el TSP.\cite{de14} Esto no prueba que no exista, pero sí significa que todo lo anterior es transferencia desde problemas parecidos.

Dificultades comunes. Se trata de observaciones propuestas por este trabajo:
\begin{itemize}[leftmargin=*]
\item Mesetas. $b$ solo toma valores enteros, por lo que muchas rutas empatan. Con selección estricta, una candidata igual se descarta y la población puede quedar inmóvil.
\item Grafos dispersos. Una ruta al azar tiene en promedio $2m/(n-1)$ conexiones válidas, donde $m$ es el número de aristas. En un grafo de 3000 vértices y grado 3 serían unas 3 de 3000, así que la búsqueda arranca casi sin información útil.
\item Redundancia. Dos escrituras del mismo ciclo pueden estar lejos como vectores, lo que debilita el sentido de la diferencia.
\item Calidad engañosa. Una ruta con una conexión faltante no necesariamente está cerca de un ciclo; quizá ningún cambio pequeño la repare.
\end{itemize}

\subsection{Criterios para interpretar un resultado}
Un ciclo producido por el método se revisa con un procedimiento determinista: contar vértices, comprobar que no se repiten, consultar cada pareja consecutiva y consultar la arista de cierre. Esa revisión sí confirma un resultado positivo. En cambio, si se agota el presupuesto, la conclusión correcta es que esta ejecución no encontró un ciclo. En el FHCP se sabe además que el ciclo existe,\cite{de3} de modo que un fallo informa sobre el método y no sobre el grafo.

Para comparar una futura implementación con otros métodos del equipo de manera justa, todos tendrían que resolver las mismas instancias, usar el mismo verificador y contar el esfuerzo igual. También convendría guardar la semilla de cada ejecución y usar los mismos puntos de partida al comparar.\cite{de5} Este documento no reporta tiempos ni tasas de éxito reales, porque esas cifras exigirían ejecuciones y un protocolo de medición; todos los números de sus ejemplos son didácticos.

\section{Dudas para el profesor}\label{de:sec:dudas}
\begin{enumerate}[leftmargin=*]
\item ¿Precisión, recall y F1 son obligatorios o basta justificar por qué no aplican? Si son obligatorios, ¿qué definición de positivo y qué referencia se aceptan?
\item Si existen varios ciclos válidos, ¿cómo debe tratarse la referencia?
\item ¿El éxito es únicamente un ciclo completo, o también se valoran rutas parciales?
\item ¿Cuántas ejecuciones por instancia, qué presupuesto y qué regla para fallos se esperan?
\item ¿Qué debe reportarse como costo computacional: evaluaciones, tiempo o ambos?
\item ¿Se añadirán grafos dirigidos fuera del conjunto FHCP? La integración considera grafos simples no dirigidos para este benchmark.
\item ¿Se permite combinar DE con búsqueda local, o debe estudiarse sola?
\item ¿Qué instancias y tamaños del FHCP se usarán, si es el caso?
\end{enumerate}

